% Ad Atticum 4.12 --- headnote
% ad-atticum-04-12, May 55 BC, Rome

\textit{Cicero to Atticus, written at Rome in May 55 BC --- one
of the shortest of the surviving Atticus letters, a single
section of practical arrangements. The Halimetus affair was an
issue between Egnatius and the jurist C. Aquilius Gallus on
which Cicero had pressed Egnatius at Antium; Macro is one of
Atticus's people, an auction at Larinum on the Ides will keep
Cicero from him; the Kalends and the day after are to be a
dinner sequence (the Crassipes gardens then home, then dinner
chez Cicero with Pilia) so he may be at Milo's door first
thing.}

\textit{The line ``facio fraudem senatus consulto'' --- ``I am
playing a trick on the senate's decree'' --- is the letter's
single moment of light. The senate's decree in question is
probably the sumptuary regulation governing dinners at home;
the Crassipes gardens lay just outside the city-line, and a
dinner there counted as one ``at an inn'' rather than ``at
home,'' beneath the decree's reach. The whole letter is a
flicker of late-spring 55 BC normality between the heavier
political letters of the year.}
